% ==================================================
% DEFINIȚII STILURI GLOBALE (TIKZ)
% ==================================================
\tikzset{
    % Stil MUX: Text centrat, în interior
    mux/.style={
        draw, rectangle, minimum width=2.0cm, minimum height=3.5cm, thick, fill=white, 
        align=center, font=\small\bfseries
    },
    % Stil CBB: Text centrat, în interior
    cbb/.style={
        draw, rectangle, minimum width=2.5cm, minimum height=2.5cm, thick, fill=white, 
        align=center, font=\bfseries
    },
    wire/.style={thick, ->, >=latex},
    % Stil pentru eticheta binară (00, 01...): La stânga pinului, font mic
    binlbl/.style={font=\tiny, anchor=east, xshift=-1pt, inner sep=1pt}
}
\newcommand{\tablestretch}{1.4}
\newcommand{\tablefontsize}{\footnotesize}

% --- Q3 ---
\subsection{Implementarea etajului $Q_3$}
\textbf{Selectori:} $Q_1, Q_0$ (Liniile hărții). Logica intrărilor $D_i$ depinde de $Q_2$ (Coloane).

\begin{figure}[H]
    \centering
    \begin{tikzpicture}[font=\footnotesize]
        % --- TABELE SIDE-BY-SIDE ---
        % Tabel J (Stânga)
        \node (tableJ) at (0,0) [anchor=north west] {
            \renewcommand{\arraystretch}{\tablestretch}
            \tablefontsize
            \begin{tabular}{|c|c|c|}
                \hline
                \rowcolor{green!20}
                \diagbox[width=2.2cm]{$Q_1 Q_0$}{$Q_3 Q_2$} & \textbf{00} & \textbf{01} \\ \hline
                \textbf{00} ($D_0$) & $\bar{P}$ & $0$ \\ \hline
                \textbf{01} ($D_1$) & $0$ & $0$ \\ \hline
                \textbf{11} ($D_3$) & $0$ & $0$ \\ \hline
                \textbf{10} ($D_2$) & $0$ & $0$ \\ \hline
            \end{tabular}
        };
        \node[above, green!50!black] at (tableJ.north) {\textbf{Zona J ($Q_3=0$)}};

        % Tabel K (Dreapta - Spațiat)
        \node (tableK) at (7.5,0) [anchor=north west] {
            \renewcommand{\arraystretch}{\tablestretch}
            \tablefontsize
            \begin{tabular}{|c|c|c|}
                \hline
                \rowcolor{red!20}
                \diagbox[width=2.2cm]{$Q_1 Q_0$}{$Q_3 Q_2$} & \textbf{11} & \textbf{10} \\ \hline
                \textbf{00} ($D_0$) & $0$ & $0$ \\ \hline
                \textbf{01} ($D_1$) & $0$ & $0$ \\ \hline
                \textbf{11} ($D_3$) & $1$ & $0$ \\ \hline
                \textbf{10} ($D_2$) & $0$ & $\bar{R}\bar{T}$ \\ \hline
            \end{tabular}
        };
        \node[above, red!50!black] at (tableK.north) {\textbf{Zona K ($Q_3=1$)}};

        % --- SCHEMA LOGICĂ ---
        \def\topY{-6.0}     
        \def\botY{-10.5}    
        \def\midY{-8.25}    
        \def\muxX{3.0}
        \def\cbbX{10.0}
        \def\pinOffset{0.6}

        % Componente cu Text ÎNĂUNTRU
        \node[mux] (muxJ) at (\muxX, \topY) {MUX-J\\4:1};
        \node[mux] (muxK) at (\muxX, \botY) {MUX-K\\4:1};
        \node[cbb] (jk) at (\cbbX, \midY) {CBB JK};

        % Conexiuni MUX -> CBB (90 grade)
        \draw[wire] (muxJ.east) -- ++(1.5,0) |- ($(jk.west)+(0,\pinOffset)$) node[pos=0.7, above] {$J_3$};
        \draw[wire] (muxK.east) -- ++(1.5,0) |- ($(jk.west)+(0,-\pinOffset)$) node[pos=0.7, above] {$K_3$};

        % Intrări MUX J
        \def\muxPinSpacing{0.65}
        \foreach \i/\lbl/\val in {
          1.5/00/$\overline{Q_2} \cdot \bar{P}$,
            0.5/01/$0$,
            -0.5/10/$0$,
            -1.5/11/$0$
        } {
            % Linia se desenează spre stânga (-0.8). Eticheta este pusă la mijloc (midway) și deasupra (above).
            \draw (muxJ.west) ++(0, \i*\muxPinSpacing) -- node[midway, above, font=\tiny, inner sep=1pt] {\lbl} ++(-0.8, 0) coordinate(startline);
            \node[left, align=right] at (startline) {\small \val};
        }

        % Intrări MUX K
        \foreach \i/\lbl/\val in {
            1.5/00/$0$,
            0.5/01/$0$,
            -0.5/10/$\overline{Q_2} \cdot \bar{R}\bar{T}$,
            -1.5/11/$Q_2$
        } {
            \draw (muxK.west) ++(0, \i*\muxPinSpacing) -- node[midway, above, font=\tiny, inner sep=1pt] {\lbl} ++(-0.8, 0) coordinate(startline);
            \node[left, align=right] at (startline) {\small \val};
        }
        
        % Ieșiri CBB
        \draw[wire] ($(jk.east)+(0,0.5)$) -- ++(1,0) node[right] {$Q_3$};
        \draw[wire] ($(jk.east)+(0,-0.5)$) -- ++(1,0) node[right] {$\overline{Q_3}$};

        % Selectori
        \draw[thick] (muxJ.south) -- (muxK.north);
        \draw[thick, <-] (muxK.south) -- ++(0, -0.5) node[below] {$Q_1, Q_0$};

    \end{tikzpicture}
    \caption{Implementarea grafică a etajului $Q_3$}
\end{figure}

\newpage

% --- Q2 ---
\subsection{Implementarea etajului $Q_2$}
\textbf{Selectori:} $Q_1, Q_0$ (Liniile hărții). Logica intrărilor $D_i$ depinde de $Q_3$ (Coloane).

\begin{figure}[H]
    \centering
    \begin{tikzpicture}[font=\footnotesize]
        % Tabel J
        \node (tableJ) at (0,0) [anchor=north west] {
            \renewcommand{\arraystretch}{\tablestretch}
            \tablefontsize
            \begin{tabular}{|c|c|c|}
                \hline
                \rowcolor{green!20}
                \diagbox[width=2.2cm]{$Q_1 Q_0$}{$Q_3 Q_2$} & \textbf{00} & \textbf{10} \\ \hline
                \textbf{00} ($D_0$) & $0$ & $0$ \\ \hline
                \textbf{01} ($D_1$) & $0$ & $0$ \\ \hline
                \textbf{11} ($D_3$) & $R$ & $0$ \\ \hline
                \textbf{10} ($D_2$) & $R$ & $\bar{R}T$ \\ \hline
            \end{tabular}
        };
        \node[above, green!50!black] at (tableJ.north) {\textbf{Zona J ($Q_2=0$)}};

        % Tabel K
        \node (tableK) at (7.5,0) [anchor=north west] {
            \renewcommand{\arraystretch}{\tablestretch}
            \tablefontsize
            \begin{tabular}{|c|c|c|}
                \hline
                \rowcolor{red!20}
                \diagbox[width=2.2cm]{$Q_1 Q_0$}{$Q_3 Q_2$} & \textbf{01} & \textbf{11} \\ \hline
                \textbf{00} ($D_0$) & $0$ & $0$ \\ \hline
                \textbf{01} ($D_1$) & $0$ & $0$ \\ \hline
                \textbf{11} ($D_3$) & $1$ & $1$ \\ \hline
                \textbf{10} ($D_2$) & $1$ & $0$ \\ \hline
            \end{tabular}
        };
        \node[above, red!50!black] at (tableK.north) {\textbf{Zona K ($Q_2=1$)}};

        % Parametri
        \def\topY{-6.0} \def\botY{-10.5} \def\midY{-8.25}
        \def\muxX{3.0} \def\cbbX{10.0} \def\pinOffset{0.6} \def\muxPinSpacing{0.65}

        % Componente
        \node[mux] (muxJ) at (\muxX, \topY) {MUX-J\\4:1};
        \node[mux] (muxK) at (\muxX, \botY) {MUX-K\\4:1};
        \node[cbb] (jk) at (\cbbX, \midY) {CBB JK};

        % Conexiuni
        \draw[wire] (muxJ.east) -- ++(1.5,0) |- ($(jk.west)+(0,\pinOffset)$) node[pos=0.7, above] {$J_2$};
        \draw[wire] (muxK.east) -- ++(1.5,0) |- ($(jk.west)+(0,-\pinOffset)$) node[pos=0.7, above] {$K_2$};

        % Intrări
        \foreach \i/\lbl/\val in {
            1.5/00/$0$,
            0.5/01/$0$,
            -0.5/10/$\overline{Q_3} \cdot R + Q_3 \cdot \bar{R}T$,
            -1.5/11/$\overline{Q_3} \cdot R$
        } {
            \draw (muxJ.west) ++(0, \i*\muxPinSpacing) -- node[midway, above, font=\tiny, inner sep=1pt] {\lbl} ++(-0.8, 0) coordinate(startline);
            \node[left, align=right] at (startline) {\small \val};
        }
        \foreach \i/\lbl/\val in {
            1.5/00/$0$,
            0.5/01/$0$,
            -0.5/10/$\overline{Q_3}$,
            -1.5/11/$1$
        } {
            \draw (muxK.west) ++(0, \i*\muxPinSpacing) -- node[midway, above, font=\tiny, inner sep=1pt] {\lbl} ++(-0.8, 0) coordinate(startline);
            \node[left, align=right] at (startline) {\small \val};
        }

        \draw[wire] ($(jk.east)+(0,0.5)$) -- ++(1,0) node[right] {$Q_2$};
        \draw[wire] ($(jk.east)+(0,-0.5)$) -- ++(1,0) node[right] {$\overline{Q_2}$};
        \draw[thick] (muxJ.south) -- (muxK.north);
        \draw[thick, <-] (muxK.south) -- ++(0, -0.5) node[below] {$Q_1, Q_0$};
    \end{tikzpicture}
    \caption{Implementarea grafică a etajului $Q_2$}
\end{figure}

\newpage

% --- Q1 ---
\subsection{Implementarea etajului $Q_1$}
\textbf{Selectori:} $Q_3, Q_2$ (Coloanele hărții). Logica intrărilor $D_i$ depinde de $Q_0$ (Linii).

\begin{figure}[H]
    \centering
    \begin{tikzpicture}[font=\footnotesize]
        % --- TABELE SIDE-BY-SIDE (Late) ---
        % Tabel J (Stânga)
        \node (tableJ) at (0,0) [anchor=north west] {
            \renewcommand{\arraystretch}{\tablestretch}
            \tablefontsize
            \setlength{\tabcolsep}{3pt}
            \begin{tabular}{|c|c|c|c|c|}
                \hline
                \rowcolor{green!20}
                \diagbox[width=2.4cm]{$Q_1 Q_0$}{$Q_3 Q_2$} & \textbf{00} & \textbf{01} & \textbf{11} & \textbf{10} \\ \hline
                \textbf{00} & $\bar{P}\bar{C}$ & $\bar{H}$ & $0$ & $1$ \\ \hline
                \textbf{01} & $1$ & $1$ & $\bar{T}$ & $1$ \\ \hline
            \end{tabular}
        };
        \node[above, green!50!black] at (tableJ.north) {\textbf{Zona J ($Q_1=0$)}};

        % Tabel K (Dreapta - Lângă J)
        % Ajustat poziția X la 8.2 pentru a face loc tabelelor mai mari
        \node (tableK) at (8.2,0) [anchor=north west] {
            \renewcommand{\arraystretch}{\tablestretch}
            \tablefontsize
            \setlength{\tabcolsep}{3pt}
            % MĂRIT WIDTH la 2.4cm
            \begin{tabular}{|c|c|c|c|c|}
                \hline
                \rowcolor{red!20}
                \diagbox[width=2.4cm]{$Q_1 Q_0$}{$Q_3 Q_2$} & \textbf{00} & \textbf{01} & \textbf{11} & \textbf{10} \\ \hline
                \textbf{11} & $R$ & $1$ & $1$ & $0$ \\ \hline
                \textbf{10} & $R$ & $1$ & $0$ & $\bar{R}$ \\ \hline
            \end{tabular}
        };
        \node[above, red!50!black] at (tableK.north) {\textbf{Zona K ($Q_1=1$)}};

        % Parametri
        \def\topY{-6.0} \def\botY{-10.5} \def\midY{-8.25}
        \def\muxX{3.5} \def\cbbX{10.5} \def\pinOffset{0.6} \def\muxPinSpacing{0.65}

        % Componente
        \node[mux] (muxJ) at (\muxX, \topY) {MUX-J\\4:1};
        \node[mux] (muxK) at (\muxX, \botY) {MUX-K\\4:1};
        \node[cbb] (jk) at (\cbbX, \midY) {CBB JK};

        % Conexiuni
        \draw[wire] (muxJ.east) -- ++(1.5,0) |- ($(jk.west)+(0,\pinOffset)$) node[pos=0.7, above] {$J_1$};
        \draw[wire] (muxK.east) -- ++(1.5,0) |- ($(jk.west)+(0,-\pinOffset)$) node[pos=0.7, above] {$K_1$};

        % Intrări
        \foreach \i/\lbl/\val in {
            1.5/00/$Q_0 + \bar{P}\bar{C}$,
            0.5/01/$Q_0 + \bar{H}$,
            -0.5/10/$1$, 
            -1.5/11/$Q_0 \cdot \bar{T}$
        } {
            \draw (muxJ.west) ++(0, \i*\muxPinSpacing) -- node[midway, above, font=\tiny, inner sep=1pt] {\lbl} ++(-0.8, 0) coordinate(startline);
            \node[left, align=right] at (startline) {\small \val};
        }
        \foreach \i/\lbl/\val in {
            1.5/00/$R$,
            0.5/01/$1$, 
            -0.5/10/$\overline{Q_0} \cdot \bar{R}$, 
            -1.5/11/$Q_0$
        } {
            \draw (muxK.west) ++(0, \i*\muxPinSpacing) -- node[midway, above, font=\tiny, inner sep=1pt] {\lbl} ++(-0.8, 0) coordinate(startline);
            \node[left, align=right] at (startline) {\small \val};
        }

        \draw[wire] ($(jk.east)+(0,0.5)$) -- ++(1,0) node[right] {$Q_1$};
        \draw[wire] ($(jk.east)+(0,-0.5)$) -- ++(1,0) node[right] {$\overline{Q_1}$};
        \draw[thick] (muxJ.south) -- (muxK.north);
        \draw[thick, <-] (muxK.south) -- ++(0, -0.5) node[below] {$Q_3, Q_2$};
    \end{tikzpicture}
    \caption{Implementarea grafică a etajului $Q_1$}
\end{figure}

\newpage

% --- Q0 ---
\subsection{Implementarea etajului $Q_0$}
\textbf{Selectori:} $Q_3, Q_2$ (Coloanele hărții). Logica intrărilor $D_i$ depinde de $Q_1$ (Linii).

\begin{figure}[H]
    \centering
    \begin{tikzpicture}[font=\footnotesize]
        % --- TABELE SIDE-BY-SIDE (Late) ---
        % Tabel J
        \node (tableJ) at (0,0) [anchor=north west] {
            \renewcommand{\arraystretch}{\tablestretch}
            \tablefontsize
            \setlength{\tabcolsep}{3pt}
            \begin{tabular}{|c|c|c|c|c|}
                \hline
                \rowcolor{green!20}
                \diagbox[width=2.4cm]{$Q_1 Q_0$}{$Q_3 Q_2$} & \textbf{00} & \textbf{01} & \textbf{11} & \textbf{10} \\ \hline
                \textbf{00} & $P$ & $1$ & $1$ & $0$ \\ \hline
                \textbf{10} & $\bar{R}$ & $0$ & $1$ & $R$ \\ \hline
            \end{tabular}
        };
        \node[above, green!50!black] at (tableJ.north) {\textbf{Zona J ($Q_0=0$)}};

        % Tabel K
        \node (tableK) at (8.2,0) [anchor=north west] {
            \renewcommand{\arraystretch}{\tablestretch}
            \tablefontsize
            \setlength{\tabcolsep}{3pt}
            % MĂRIT WIDTH la 2.4cm
            \begin{tabular}{|c|c|c|c|c|}
                \hline
                \rowcolor{red!20}
                \diagbox[width=2.4cm]{$Q_1 Q_0$}{$Q_3 Q_2$} & \textbf{00} & \textbf{01} & \textbf{11} & \textbf{10} \\ \hline
                \textbf{01} & $1$ & $1$ & $\bar{T}$ & $1$ \\ \hline
                \textbf{11} & $R$ & $1$ & $1$ & $1$ \\ \hline
            \end{tabular}
        };
        \node[above, red!50!black] at (tableK.north) {\textbf{Zona K ($Q_0=1$)}};

        % Parametri
        \def\topY{-6.0} \def\botY{-10.5} \def\midY{-8.25}
        \def\muxX{3.5} \def\cbbX{10.5} \def\pinOffset{0.6} \def\muxPinSpacing{0.65}

        % Componente
        \node[mux] (muxJ) at (\muxX, \topY) {MUX-J\\4:1};
        \node[mux] (muxK) at (\muxX, \botY) {MUX-K\\4:1};
        \node[cbb] (jk) at (\cbbX, \midY) {CBB JK};

        % Conexiuni
        \draw[wire] (muxJ.east) -- ++(1.5,0) |- ($(jk.west)+(0,\pinOffset)$) node[pos=0.7, above] {$J_0$};
        \draw[wire] (muxK.east) -- ++(1.5,0) |- ($(jk.west)+(0,-\pinOffset)$) node[pos=0.7, above] {$K_0$};

        % Intrări
        \foreach \i/\lbl/\val in {
          1.5/00/$\overline{Q_1} \cdot P +Q_1 \cdot \bar{R}$, 
            0.5/01/$\overline{Q_1}$, 
            -0.5/10/$Q_1 \cdot R$, 
            -1.5/11/$1$
        } {
            \draw (muxJ.west) ++(0, \i*\muxPinSpacing) -- node[midway, above, font=\tiny, inner sep=1pt] {\lbl} ++(-0.8, 0) coordinate(startline);
            \node[left, align=right] at (startline) {\small \val};
        }
        \foreach \i/\lbl/\val in {
            1.5/00/$Q_1 + R$, 
            0.5/01/$1$, 
            -0.5/10/$1$, 
            -1.5/11/$\overline{Q_1}+\bar{T}$
        } {
            \draw (muxK.west) ++(0, \i*\muxPinSpacing) -- node[midway, above, font=\tiny, inner sep=1pt] {\lbl} ++(-0.8, 0) coordinate(startline);
            \node[left, align=right] at (startline) {\small \val};
        }

        \draw[wire] ($(jk.east)+(0,0.5)$) -- ++(1,0) node[right] {$Q_0$};
        \draw[wire] ($(jk.east)+(0,-0.5)$) -- ++(1,0) node[right] {$\overline{Q_0}$};
        \draw[thick] (muxJ.south) -- (muxK.north);
        \draw[thick, <-] (muxK.south) -- ++(0, -0.5) node[below] {$Q_3, Q_2$};
    \end{tikzpicture}
    \caption{Implementarea grafică a etajului $Q_0$}
\end{figure}
